\documentclass[11pt,a4paper]{article}
\usepackage{fontspec}
\setmainfont{DejaVu Sans}[Scale=0.90]
\usepackage{amsmath,amssymb}
\usepackage[margin=2.4cm]{geometry}
\usepackage{booktabs,array,longtable,tabularx}
\usepackage{enumitem}
\usepackage{titlesec}
\usepackage{parskip}
\titlelabel{\thetitle.\quad}
\titleformat{\section}{\normalfont\large\bfseries}{\thesection.}{0.6em}{}
\titleformat{\subsection}{\normalfont\normalsize\bfseries}{\thesubsection}{0.6em}{}
\titlespacing*{\section}{0pt}{16pt}{6pt}
\titlespacing*{\subsection}{0pt}{12pt}{4pt}
\setlist[itemize]{leftmargin=1.4em,itemsep=1pt,topsep=3pt}
\newcolumntype{L}[1]{>{\raggedright\arraybackslash}p{#1}}
\begin{document}
\begin{center}
{\Large\bfseries Fiche de lecture --- Krapez (2018)}\\[4pt]
{\small Krapez, \emph{International Journal of Thermal Sciences} \textbf{136} (2018) 182--199}\\[2pt]
{\small Version corrigée --- 8 septembre 2026}
\end{center}

space{4pt}


\section{Identification}

\begin{center}
\begin{tabularx}{\textwidth}{@{}lX@{}}
\toprule
Auteur & Jean-Claude Krapez, ONERA, département DOTA, Salon-de-Provence \\
Titre & \emph{Linear, trigonometric and hyperbolic profiles of thermal effusivity in the Liouville
space and related quadrupoles: Simple analytical tools for modeling graded layers and multilayers} \\
Revue & International Journal of Thermal Sciences, volume 136, pages 182--199, 2018 \\
DOI & 10.1016/j.ijthermalsci.2018.10.018 \\
Accès & Postprint en libre accès, dépôt HAL, identifiant hal-01974990 \\
\bottomrule
\end{tabularx}
\end{center}

\section{Nature et objectif}

Article théorique et méthodologique, destiné à la communauté de la métrologie thermique et de la
modélisation des matériaux à gradient de propriétés. Il ne comporte aucune donnée expérimentale.

L'objectif annoncé est la construction de nouvelles familles de profils de propriétés pour lesquels
l'équation de la chaleur unidimensionnelle admet une solution analytique exacte, exprimée uniquement
au moyen de fonctions élémentaires. Les solutions antérieures pour milieux gradués font intervenir des
fonctions spéciales --- Bessel modifiées, Airy, Kelvin --- coûteuses à évaluer et peu maniables.

\section{Résultat principal}

Après transformation de Liouville, l'équation de la chaleur en milieu à gradient prend la forme normale
\[ \frac{d^{2}\psi}{d\xi^{2}} - \bigl(V(\xi) + p\bigr)\psi = 0 ,
\qquad V(\xi) = \frac{s''(\xi)}{s(\xi)} , \qquad s = b^{\pm 1/2} , \]
où $\xi$ désigne la racine du temps de diffusion intégré, $b$ l'effusivité thermique et $s$ la
métapropriété. L'imposition d'un potentiel constant, $V = \beta$, engendre exactement trois familles de
profils solvables, auxquelles sont associées des matrices quadripolaires ne faisant intervenir que des
fonctions hyperboliques et rationnelles.

\section{Structure de l'article}

\begin{center}
\begin{tabularx}{\textwidth}{@{}lX@{}}
\toprule
Section 1 & Introduction et état de l'art des profils solvables antérieurs \\
Section 2 & Résolution dans l'espace de Liouville : les deux équations duales, la transformation, la
prééminence de l'effusivité, la méthode des quadripôles, l'indétermination de la transformation
inverse \\
Section 3 & Les trois familles issues d'un potentiel constant : profils, retour à l'espace physique,
quadripôles, réponses photothermiques calculées \\
Section 4 & Discussion, hiérarchie des modèles multicouches, transpositions à d'autres domaines \\
Appendice & Construction de la matrice quadripolaire à partir des solutions indépendantes \\
\bottomrule
\end{tabularx}
\end{center}

\section{Hypothèses de l'article}

\begin{center}
\begin{tabularx}{\textwidth}{@{}lXX@{}}
\toprule
\textbf{Élément} & \textbf{Hypothèse} & \textbf{Portée} \\
\midrule
Géométrie & Transfert unidimensionnel & Restriction explicitement soulignée par l'auteur \\
\addlinespace
Conduction & Loi de Fourier & Le cadre non-Fourier n'est pas abordé \\
\addlinespace
Propriétés & Gradient continu ; les couches homogènes en sont un cas particulier &
Applicable aux profils gradués comme aux empilements \\
\addlinespace
Régime & Harmonique et transitoire & Les deux formulations sont établies \\
\addlinespace
Source & Flux surfacique uniforme & Les sources volumiques ne sont pas traitées \\
\addlinespace
Face arrière & Substrat semi-infini dans la configuration de référence & Toute autre condition
requiert une modification de l'impédance de fermeture \\
\addlinespace
Interfaces & Contact thermique parfait & Une résistance de contact nécessite un quadripôle
supplémentaire \\
\addlinespace
Condition initiale & Nulle dans la formulation transitoire & Requise par la transformée employée \\
\addlinespace
Température & Propriétés indépendantes de $T$ & Les fortes non-linéarités et les changements de phase
sortent du cadre \\
\bottomrule
\end{tabularx}
\end{center}

\section{Compatibilité avec le présent travail}

Les hypothèses de géométrie unidimensionnelle, de contact parfait et de propriétés indépendantes de la
température sont communes aux deux cadres. La divergence porte sur la loi de conduction : l'article se
place sous la loi de Fourier, tandis que le présent travail examine le cas d'une loi à temps de
relaxation fini.

Cette divergence est structurante. Sous la loi de Fourier, le paramètre spectral vaut $p = i\omega$ et
l'énergie associée à la forme de Schrödinger demeure sur l'axe imaginaire. Sous une loi de type
Cattaneo, ce paramètre devient $p + \tau p^{2}$ et l'énergie acquiert une partie réelle.

\section{Suite de lecture}

La lecture en seconde passe porte sur les sections 2.2 à 2.5, qui contiennent la totalité de la
machinerie théorique. Les sections 3 et 4 déroulent l'application de cette machinerie aux trois
familles de profils et n'introduisent aucun élément conceptuel nouveau.

La référence [58] de l'article --- Krapez, \emph{International Journal of Heat and Mass Transfer}
\textbf{99} (2016) 485--503 --- constitue le travail fondateur dont le présent article est la suite.
Elle expose la méthode PROFIDT et les transformations de Darboux appliquées conjointement au profil et
au champ.
\end{document}